\documentclass[10pt]{article}
\usepackage[T1]{fontenc}
\usepackage[utf8]{inputenc}
\usepackage[french]{babel}
\usepackage{amsmath,amssymb,mathrsfs,mathpazo}
\usepackage{geometry}
\geometry{a4paper,margin=2cm}
\usepackage{graphicx,array,tabularx,tabularray,colortbl,multirow,multicol,booktabs,enumitem,pifont,ulem,fancybox,diagbox,cancel}
\usepackage{pstricks,pstricks-add,pst-plot,pst-tree,pst-node,pst-text,pst-func,pst-math,pst-bspline,pst-3dplot,pst-eucl,pst-solides3d}
\let\pscalculate\undefined
\usepackage{pst-bezier}
\usepackage{tikz,pgfplots,tkz-tab}
\usetikzlibrary{arrows,arrows.meta,patterns,calc,positioning,trees,decorations.markings,shapes.geometric,shapes.misc,shapes.arrows,intersections,decorations.pathmorphing,decorations.pathreplacing,backgrounds,shadows,babel}
\pgfplotsset{compat=1.18}
\usepgfplotslibrary{fillbetween}
\usepackage[np]{numprint}
\usepackage{listings,makecell,pas-tableur,twoopt}
\let\barre\undefined
\usepackage[pstricks]{tablvar}
\definecolor{darkgreen}{RGB}{0,100,0}
\providecommand{\red}{\color{red}}
\providecommand{\blue}{\color{blue}}
\providecommand{\green}{\color{green}}
\providecommand{\black}{\color{black}}
\providecommand{\white}{\color{white}}
\providecommand{\gray}{\color{gray}}
\providecommand{\cyan}{\color{cyan}}
\providecommand{\magenta}{\color{magenta}}
\providecommand{\yellow}{\color{yellow}}
\providecommand{\violet}{\color{violet}}
\providecommand{\orange}{\color{orange}}
\providecommand{\euro}{\texteuro}
\providecommand{\EUR}{\texteuro}
\providecommand{\ieme}{\textsuperscript{e}}
\providecommand{\degres}{\ensuremath{^\circ}}
\providecommand{\degree}{\ensuremath{^\circ}}
\providecommand{\vv}[1]{\overrightarrow{#1}}
\providecommand{\Coord}[1]{#1}
\providecommand{\decosix}{}
\providecommand{\starredbullet}{\ensuremath{\bullet}}
\pagestyle{empty}
\setlength{\parindent}{0pt}
\usepackage{adjustbox}\DeclareUnicodeCharacter{2236}{:}\DeclareUnicodeCharacter{2190}{\ensuremath{\leftarrow}}
\providecommand{\R}{\mathbb{R}}
\providecommand{\N}{\mathbb{N}}
\providecommand{\Z}{\mathbb{Z}}
\providecommand{\Q}{\mathbb{Q}}
\providecommand{\D}{\mathbb{D}}
\providecommand{\C}{\mathbb{C}}
\providecommand{\e}{{\mathrm{e}}}
\providecommand{\vect}[1]{\overrightarrow{#1}}
\providecommand{\vectt}[1]{\overrightarrow{\mathrm{#1}}}
\providecommand{\vvt}[1]{\overrightarrow{\mathrm{#1}}}
\providecommand{\np}[2][]{{#2~#1}}
\providecommand{\vv}[1]{{\overrightarrow{#1}}}
\renewcommand{\d}{{\mathrm{d}}}
\providecommand{\degres}{{^\circ}}
\providecommand{\degre}{{^\circ}}
\providecommand{\blue}{\color{blue}}
\providecommand{\red}{\color{red}}
\providecommand{\enskip}{\quad}
\providecommand{\llbracket}{\lbrack\!\lbrack}
\providecommand{\rrbracket}{\rbrack\!\rbrack}
\providecommand{\euro}{{\text{€}}}
\providecommand{\no}{{\text{n°}}}
\providecommand{\interval}[2]{{[#1;#2]}}
\providecommand{\intervFO}[2]{{[#1;#2[}}
\providecommand{\dotfill}{\ldots}

\providecommand{\ds}{\displaystyle}
\providecommand{\textding}[1]{#1}
\providecommand{\Oij}{\(\mathrm{O};\vec{\imath},\vec{\jmath}\)}
\providecommand{\Oijk}{\(\mathrm{O};\vec{\imath},\vec{\jmath},\vec{k}\)}
\providecommand{\Ouv}{\(\mathrm{O};\vec{u},\vec{v}\)}
\providecommand{\intervFF}[2]{\ensuremath{[#1;#2]}}
\renewcommand{\euro}{\ifmmode\text{\texteuro}\else\texteuro\fi}
\DeclareUnicodeCharacter{00B0}{\ifmmode^{\circ}\else\textdegree\fi}
\DeclareUnicodeCharacter{20AC}{\ifmmode\text{\texteuro}\else\texteuro\fi}
% Énoncés : Éducation nationale. Transcription et corrigés : APMEP, Denis Vergès et contributeurs. Découpage et classement : CoopMaths, Sébastien Lozano, Éric Elter et contributeurs. Adaptation de la présentation : Maths Explorer.
% CC BY-SA : https://coopmaths.fr/www/about/
% https://coopmaths.fr/alea/static/bac/2022/tex/bac_2022_05_sujet2_asie_3.tex
\begin{document}
\noindent\textit{APMEP / CoopMaths. CC BY-SA. Présentation adaptée par Maths Explorer.}\par\medskip
\medskip

\emph{Principaux domaines abordés}: Probabilités conditionnelles et indépendance. Variables aléatoires.

\medskip

\begin{center}\emph{Les deux parties de cet exercice sont indépendantes.}\end{center}

\textbf{Partie 1}

\medskip

Julien doit prendre l'avion; il a prévu de prendre le bus pour se rendre à l'aéroport. 

S'il prend le bus de $8$~h, il est sûr d'être à l'aéroport à temps pour son vol.

Par contre, le bus suivant ne lui permettrait pas d'arriver à temps à l'aéroport.

Julien est parti en retard de son appartement et la probabilité qu'il manque son bus est de $0,8$.

S'il manque son bus, il se rend à l'aéroport en prenant une compagnie de voitures privées; il a alors une probabilité de $0,5$ d'être à l'heure à l'aéroport.

\smallskip

On notera :

\setlength\parindent{12mm}
\begin{itemize}
\item[$\bullet~~$] $B$ l'évènement: \og Julien réussit à prendre son bus \fg ;
\item[$\bullet~~$] $V$ l'évènement: \og Julien est à l'heure à l'aéroport pour son vol \fg.
\end{itemize}
\setlength\parindent{0mm}

\medskip

\begin{enumerate}
\item Donner la valeur de $P_B(V)$.
\item Représenter la situation par un arbre pondéré.
\item Montrer que $P(V) = 0,6$.
\item Si Julien est à l'heure à l'aéroport pour son vol, quelle est la probabilité qu'il soit arrivé à l'aéroport en bus ? Justifier.
\end{enumerate}

\bigskip

\textbf{Partie 2}

\medskip

Les compagnies aériennes vendent plus de billets qu'il n'y a de places dans les avions car certains passagers ne se présentent pas à l'embarquement du vol sur lequel ils ont réservé.
On appelle cette pratique le surbooking.

Au vu des statistiques des vols précédents, la compagnie aérienne estime que chaque passager a 5\,\% de chance de ne pas se présenter à l'embarquement.

Considérons un vol dans un avion de $200$~places pour lequel $206$~billets ont été vendus.
On suppose que la présence à l'embarquement de chaque passager est indépendante des autres passagers et on appelle $X$ la variable aléatoire qui compte le nombre de passagers se présentant à l'embarquement.

\medskip

\begin{enumerate}
\item Justifier que $X$ suit une loi binomiale dont on précisera les paramètres.
\item En moyenne, combien de passagers vont-ils se présenter à l'embarquement ?
\item Calculer la probabilité que $201$ passagers se présentent à l'embarquement. Le résultat sera arrondi à $10^{-3}$ près.
\item Calculer $P(X \leqslant 200)$, le résultat sera arrondi à $10^{-3}$ près. Interpréter ce résultat dans le contexte de l'exercice.
\item La compagnie aérienne vend chaque billet à $250$ euros.

Si plus de $200$ passagers se présentent à l'embarquement, la compagnie doit rembourser le billet d'avion et payer une pénalité de $600$ euros à chaque passager lésé.

On appelle :

\setlength\parindent{12mm}
\begin{description}
\item[ ] $Y$ la variable aléatoire égale au nombre de passagers qui ne peuvent pas embarquer
bien qu'ayant acheté un billet;
\item[ ] $C$ la variable aléatoire qui totalise le chiffre d'affaire de la compagnie aérienne sur ce vol.
\end{description}

On admet que $Y$ suit la loi de probabilité donnée par le tableau suivant:

\begin{center}
\begin{tabularx}{\linewidth}{|c|*{7}{>{\centering \arraybackslash}X|}}\hline
$y_i$					&0	&1	&2	&3	&4	&5	&6\\ \hline
$P\left(Y = y_i\right)$	&\np{0,94775}&\np{0,03063}&\np{0,01441}&\np{0,00539}&\np{0,00151}&\np{0,00028}& \\ \hline
\end{tabularx}
\end{center}

	\begin{enumerate}
		\item Compléter la loi de probabilité donnée ci-dessus en calculant $P(Y = 6)$.
		\item Justifier que: $C = \np{51500} - 850Y$.
		\item Donner la loi de probabilité de la variable aléatoire $C$ sous forme d'un tableau.
		
Calculer l'espérance de la variable aléatoire $C$ à l'euro près.
		\item Comparer le chiffre d'affaires obtenu en vendant exactement $200$ billets et le chiffre d'affaires moyen obtenu en pratiquant le surbooking.
	\end{enumerate}
\end{enumerate}

\bigskip

\end{document}
